\section{Corrections and Calibrations}
\label{sec:corrections}

Corrections are applied in a fixed order: collaboration-recommended
corrections first, then analysis-specific corrections derived where a
data/MC discrepancy persists.

\subsection{Collaboration-recommended corrections}
\label{sec:corrections:recommended}
Track momentum scale and energy loss are corrected during reconstruction
(Sec.~\ref{sec:recon:fsp}).
PID efficiency and fake-rate corrections are applied offline by
\texttt{utilities.apply\_pid\_corrections()} through the systematics
framework, using per-run-period tables for the lepton, the kaon and both
pions; the per-particle weights are multiplied into a single event weight.
The correction is run dependent and is evaluated before run periods are
combined.
\AnalysisTBD{the PID table versions and their provenance; the tables are read
from absolute user-specific paths and are not part of the repository}

\subsection{Analysis-specific corrections}
\label{sec:corrections:specific}
The generic-$B\bar{B}$ reweighting of Sec.~\ref{sec:bbbar} is the main
analysis-specific correction.

\subsection{Correction ordering and double-application safeguards}
\label{sec:corrections:order}
Table~\ref{tab:correction-order} records where each correction is applied and,
critically, where it is \emph{not}.

% PROVENANCE
%   status      : CODE_FACT (application points) / ANALYSIS_DECISION (ordering)
%   produced-by : Recon_scripts/2_Reconstruction.py; utilities.py
%                 (apply_pid_corrections, create_templates_new,
%                 create_workspace); 5_BBbkg_weights_optuna_minuit.py
%   commit      : dd520a5
%   sample      : MC16rd + Proc16/Prompt16, Run 1 + Run 2, e and mu
%   selection   : n/a (applies to the whole chain)
%   corrections : the table is itself the list
\begin{table}[htbp]
  \centering
  \caption{Correction application points.  A correction must appear exactly
  once in the chain leading to any given number.}
  \label{tab:correction-order}
  \small
  \begin{tabular}{p{0.24\textwidth} p{0.30\textwidth} p{0.38\textwidth}}
    \hline
    Correction & Applied at & Note \\
    \hline
    Track momentum scale, energy loss & Reconstruction & \texttt{central} variation only \\
    PID efficiency / fake rate & Offline, before run combination & Run dependent \\
    Generic MC luminosity ($0.25$) & Normalisation comparisons & Generic MC only, not signal MC \\
    Generic-$B\bar{B}$ category weights & Sec.~\ref{sec:bbbar} & Derived in the tuning region \\
    Fake-$D$ normalisation ($0.87$ Run~1, $1.0$ Run~2) & Generic-$B\bar{B}$ tuning \emph{only} & \textbf{Not} applied when building fit templates; in the fit this normalisation is absorbed by the free parameter \texttt{bkg\_fakeD\_norm} (Sec.~\ref{sec:fit}).  Applying it at both points would double-count it \\
    Slow-$\pi^{0}$ efficiency & Not applied & Awaiting an updated MC16rd table \\
    Form factors & Not applied & See below \\
    \hline
  \end{tabular}
\end{table}

\AnalysisTBD{form-factor reweighting (e.g.\ via HAMMER) and branching-fraction
updates are not implemented in the analysis code at the pinned commit.  State
the plan, or the justification for omitting them}
